本论文是在指导教师贺华老师的悉心指导下完成的。从选题确定、技术路线设计、实验方案完善到论文撰写和修改，贺老师都给予了耐心细致的指导和严格要求。老师严谨的治学态度、清晰的工程思维和务实的工作方法，使我在毕业设计过程中不仅完成了具体系统实现，也进一步理解了科研问题分析、实验论证和论文表达之间的关系。在此向贺华老师致以诚挚的感谢。

感谢机电工程学院各位老师在本科阶段给予的培养与帮助。大学学习期间，专业课程训练为本文涉及的系统建模、程序实现、实验分析和工程表达奠定了基础，也使我能够在毕业设计中将理论学习与实际问题结合起来。感谢学院和学校提供的学习平台与科研环境，使我能够较为系统地完成本次毕业设计工作。

感谢同学和朋友在论文资料整理、实验结果核对、文档排版和答辩准备过程中给予的帮助与建议。与大家的交流使我能够及时发现论文表述和系统实现中的不足，并不断完善本文的结构、内容和实验说明。

最后，感谢家人长期以来的理解、支持与鼓励。正是他们在学习和生活中的支持，使我能够保持稳定的学习状态并顺利完成本科阶段的学习任务。谨向所有关心、帮助和支持过我的老师、同学、朋友和家人表示衷心感谢。
